% TC-0015 fragment; included by c89-test-report.tex.
\def\TCID{TC-0015}
\def\TCTitle{C89 date and time}
\def\TCPurpose{Verify all nine clock, calendar, normalization, text conversion,
and C-locale formatting functions.}
\def\TCPriority{\TCPriorityHigh}
\def\TCRequirementRef{REQ-0015}
\def\TCDesignRef{ANSI X3J11 C89 draft 4.12.}
\def\TCFeatureRef{\texttt{include/time.h}; \texttt{src/time.c};
\texttt{src/platform/windows/clock.c}; \texttt{tests/c89/time.c}.}
\def\TCPreconditions{TinyCC and Windows FILETIME APIs are available.}
\def\TCEnvironment{x64 behavioral execution, Microsoft UCRT comparison, and
Windows ARM64 cross-build.}
\def\TCAssumptions{Calendar arithmetic is proleptic Gregorian; timezone rules
come from Windows and daylight status remains indeterminate.}
\def\TCInputData{Unix epoch, current time, overflowing January date fields,
all C89 format conversions, and an undersized output buffer.}
\def\TCInitialState{Shared calendar and text buffers are unspecified.}
\def\TCProcedure{\par\begin{enumerate}
\item Verify epoch fields, weekday, year-day, and fixed asctime output.
\item Round-trip current local time and normalize January 32.
\item Exercise every C89 strftime conversion and insufficient capacity.
\item Verify wall storage and nondecreasing process time.
\item Cross-link the implementation as PE/COFF ARM64.
\end{enumerate}}
\def\TCExpectedResult{All nine functions satisfy the tested C89 behavior and
the implementation cross-builds with ARM64 machine type AA64.}
\def\TCPostConditions{Shared static conversion buffers contain the last result.}
\TCTemplate
